% ============================================================
% preamble.tex
% 2026 MathorCup 大数据竞赛 A 题论文导言区
% 编译引擎：xelatex（中文由 ctexart / xeCJK 提供）
% 说明：本文件只放宏包、全局格式与常用宏，不放正文。
% ============================================================

% ---------- 页面与版式 ----------
\usepackage{geometry}
\geometry{left=2.6cm,right=2.6cm,top=2.5cm,bottom=2.5cm}
\linespread{1.3}
\usepackage{fancyhdr}
\pagestyle{fancy}
\fancyhf{}
\fancyfoot[C]{\thepage}
\renewcommand{\headrulewidth}{0pt}
\setlength{\emergencystretch}{2em}

% ---------- 数学 ----------
\usepackage{amsmath,amssymb}
\usepackage{bm}
\allowdisplaybreaks[3]
\numberwithin{equation}{section}
\numberwithin{figure}{section}
\numberwithin{table}{section}
\DeclareMathOperator*{\argmin}{arg\,min}
\DeclareMathOperator*{\argmax}{arg\,max}
\newcommand{\R}{\mathbb{R}}
\newcommand{\E}{\mathbb{E}}
\newcommand{\Var}{\operatorname{Var}}
\newcommand{\Cov}{\operatorname{Cov}}
\newcommand{\Corr}{\operatorname{Corr}}
\newcommand{\norm}[1]{\left\lVert #1 \right\rVert}
\newcommand{\abs}[1]{\left\lvert #1 \right\rvert}
\newcommand{\inner}[2]{\left\langle #1,#2 \right\rangle}

% ---------- 一套符号：向量、矩阵、希腊字母 ----------
% 约定：粗体小写表示向量，粗体大写表示矩阵。
\newcommand{\vx}{\bm{x}}
\newcommand{\vy}{\bm{y}}
\newcommand{\bz}{\bm{z}}
\newcommand{\bw}{\bm{w}}
\newcommand{\vb}{\bm{b}}
\newcommand{\bs}{\bm{s}}
\newcommand{\bzero}{\bm{0}}
\newcommand{\bX}{\bm{X}}
\newcommand{\bY}{\bm{Y}}
\newcommand{\bZ}{\bm{Z}}
\newcommand{\bW}{\bm{W}}
\newcommand{\bB}{\bm{B}}
\newcommand{\bA}{\bm{A}}
\newcommand{\bR}{\bm{R}}
\newcommand{\bM}{\bm{M}}
\newcommand{\bI}{\bm{I}}
\newcommand{\bD}{\bm{D}}
\newcommand{\bL}{\bm{L}}
\newcommand{\bU}{\bm{U}}
\newcommand{\bV}{\bm{V}}
\newcommand{\bP}{\bm{P}}
\newcommand{\bT}{\bm{T}}
\newcommand{\bbeta}{\bm{\beta}}
\newcommand{\btheta}{\bm{\theta}}
\newcommand{\bmu}{\bm{\mu}}
\newcommand{\bSigma}{\bm{\Sigma}}
\newcommand{\bepsilon}{\bm{\epsilon}}
\newcommand{\vxt}{\tilde{\bm{x}}}
\newcommand{\yhat}{\hat{y}}
\newcommand{\Dcal}{\mathcal{D}}
\newcommand{\Lcal}{\mathcal{L}}
\newcommand{\Ncal}{\mathcal{N}}
\newcommand{\Jcal}{\mathcal{J}}
\newcommand{\Tcal}{\mathcal{T}}
\newcommand{\Mcal}{\mathcal{M}}

% ---------- 图 ----------
\usepackage{graphicx}
% 图片统一放在项目根目录的 img/ 下，由 src/ 的 Python 脚本生成。
% 同时保留 img/，便于从项目根目录编译时也能找到图。
\graphicspath{{../img/}{img/}}
\DeclareGraphicsExtensions{.pdf,.png,.jpg,.jpeg}

% ---------- 表 ----------
\usepackage{array}
\usepackage{booktabs}
\usepackage{tabularx}
\usepackage{longtable}
\renewcommand{\arraystretch}{1.18}

% ---------- 算法 ----------
\usepackage{xcolor}
\usepackage[ruled,vlined,linesnumbered]{algorithm2e}
\SetAlgorithmName{算法}{算法}{算法列表}
\SetKwInOut{Input}{输入}
\SetKwInOut{Output}{输出}

% ---------- 代码 ----------
\usepackage{listings}
\renewcommand{\lstlistingname}{代码}
\definecolor{codeblue}{RGB}{25,70,160}
\definecolor{codegreen}{RGB}{0,110,60}
\definecolor{codered}{RGB}{170,30,30}
\definecolor{codegray}{RGB}{110,110,110}
\lstdefinestyle{pycode}{
  language=Python,
  basicstyle=\ttfamily\small,
  keywordstyle=\color{codeblue}\bfseries,
  commentstyle=\color{codegreen},
  stringstyle=\color{codered},
  numbers=left,
  numberstyle=\tiny\color{codegray},
  numbersep=6pt,
  frame=single,
  breaklines=true,
  showstringspaces=false,
  columns=fullflexible,
  keepspaces=true,
  tabsize=4,
  captionpos=b
}
\lstdefinestyle{shellcode}{
  language=bash,
  basicstyle=\ttfamily\small,
  keywordstyle=\color{codeblue}\bfseries,
  commentstyle=\color{codegreen},
  stringstyle=\color{codered},
  numbers=left,
  numberstyle=\tiny\color{codegray},
  numbersep=6pt,
  frame=single,
  breaklines=true,
  showstringspaces=false,
  columns=fullflexible,
  keepspaces=true,
  tabsize=4,
  captionpos=b
}
\lstset{style=pycode}
% 注意：代码块内请使用英文注释，避免 CJK 字符在 listings 中断行异常。

% ---------- 参考文献与链接 ----------
\usepackage[numbers,sort&compress]{natbib}
\usepackage[hidelinks,bookmarksnumbered=true]{hyperref}

% ---------- 常用占位宏 ----------
% 定稿前全文搜索 TODO 与“待定”，逐项清零。
\newcommand{\todo}[1]{\textcolor{red!75!black}{\textbf{[TODO: #1]}}}
\newcommand{\tbd}{\textcolor{red!75!black}{待定}}

% \figauto{文件名}{图题}{标签}
% 用法：\figauto{fig_q1_framework.pdf}{问题一框架}{fig:q1framework}
% 若 ../img/ 或 img/ 下存在该文件则自动插图；否则显示带文件名的占位框。
% 这样骨架在图片尚未生成时也能编译，图片生成后无需改正文。
\newcommand{\figauto}[4][0.82\textwidth]{%
  \begin{figure}[htbp]
    \centering
    \IfFileExists{../img/#2}{%
      \includegraphics[width=#1]{#2}%
    }{%
      \IfFileExists{img/#2}{%
        \includegraphics[width=#1]{#2}%
      }{%
        \fbox{\begin{minipage}[c][4.2cm][c]{0.78\textwidth}
          \centering\small
          图待生成：\texttt{\detokenize{#2}}\\
          由 \texttt{src/} 下脚本输出到 \texttt{img/}
        \end{minipage}}%
      }%
    }%
    \caption{#3}
    \label{#4}
  \end{figure}%
}

% ---------- 章节标题格式 ----------
\ctexset{
  section = {format = \Large\bfseries},
  subsection = {format = \large\bfseries},
  subsubsection = {format = \normalsize\bfseries},
  paragraph = {format = \normalsize\bfseries}
}
